% Aula 10 --- uma divisão em que o Gini e a entropia discordam, e por que isso é raro.
Ana é bióloga e usa uma árvore de classificação para distinguir mudas de duas espécies, a
espécie $1$ e a espécie $0$, a partir de características medidas no viveiro. Num nó com $12$
mudas, $6$ de cada espécie, a árvore compara duas divisões:
\begin{description}[leftmargin=1.4cm,style=nextline]
  \item[A] pela altura, separa dois grupos de $6$ mudas, um com $2$ da espécie $1$ e $4$ da
        espécie $0$ e outro com $4$ da espécie $1$ e $2$ da espécie $0$;
  \item[B] por uma mancha rara na folha, isola uma única muda, da espécie $0$, das outras $11$,
        que têm $6$ da espécie $1$ e $5$ da espécie $0$.
\end{description}
Numa região com proporção $p$ de mudas da espécie $1$, o índice de Gini é $G(p)=2p(1-p)$, a
entropia é $D(p)=-p\log_2p-(1-p)\log_2(1-p)$, com $0\log_20=0$, e o erro de classificação é
$E(p)=\min(p,1-p)$. Uma divisão em regiões com $n_1$ e $n_2$ mudas e proporções $p_1$ e $p_2$
tem impureza ponderada $\big[n_1I(p_1)+n_2I(p_2)\big]/(n_1+n_2)$, para cada medida $I$, e a
árvore escolhe a divisão de menor impureza ponderada.
\begin{itens}
  \item \pts{1,4} Calcule a impureza ponderada das divisões A e B pelas três medidas e diga
        qual divisão cada medida escolhe. \emph{Dados:} $\log_23\approx1{,}5850$,
        $\log_25\approx2{,}3219$ e $\log_211\approx3{,}4594$.
  \item \pts{1,4} Seja $I$ uma medida de impureza qualquer para duas classes, com
        $I(p)=I(1-p)$, $I(0)=0$ e $I(6/11)>0$. Mostre que A tem impureza ponderada
        menor que B se e somente se $I(1/3)/I(6/11)<11/12$. Calcule essa razão para as
        três medidas e explique que propriedade da curva da entropia a leva a escolher B.
  \item \pts{1,2} No exemplo dos slides, o \texttt{DecisionTreeClassifier} é ajustado com
        \texttt{min\_samples\_leaf=5}, que só admite divisões com pelo menos $5$
        observações em cada filho. O que acontece com a divisão B? As medidas continuam
        discordando? Qual das duas divisões você preferiria, e por quê?
        As notas dizem que o Gini e a entropia quase sempre escolhem as mesmas divisões; o que
        este exemplo diz sobre esse ``quase''?
\end{itens}

\begin{solucao}
\textbf{(a)} Na divisão A, as proporções da espécie $1$ são $1/3$ e $2/3$, cada uma num filho
com metade das mudas. Na divisão B, são $0$, numa muda, e $6/11$, nas outras $11$. Pelo Gini,
\[
  \text{A: }\ \tfrac12\cdot2\cdot\tfrac13\cdot\tfrac23+\tfrac12\cdot2\cdot\tfrac23\cdot\tfrac13=\tfrac49\approx0{,}444,
  \qquad
  \text{B: }\ \tfrac1{12}\cdot0+\tfrac{11}{12}\cdot2\cdot\tfrac6{11}\cdot\tfrac5{11}=\tfrac5{11}\approx0{,}455 .
\]
Pela entropia, como $D(1/3)=D(2/3)=\tfrac13\log_23+\tfrac23\log_2\tfrac32=\log_23-\tfrac23$,
\begin{align*}
  \text{A: }\ &\log_23-\tfrac23\approx0{,}918,\\
  \text{B: }\ &\tfrac{11}{12}D\big(\tfrac6{11}\big)
  =\tfrac1{12}\Big(6\log_2\tfrac{11}6+5\log_2\tfrac{11}5\Big)
  =\tfrac1{12}\big(11\log_211-6\log_26-5\log_25\big)\approx0{,}911 ,
\end{align*}
com $\log_26=1+\log_23$. Pelo erro de classificação, A erra $2$ mudas em cada filho, $4/12=1/3$,
e B erra as $5$ mudas da espécie $0$ no filho grande, $5/12\approx0{,}417$. O Gini e o erro
escolhem A; a entropia escolhe B. O nó tem $G=0{,}5$ e $D=1$, e as reduções correspondentes são
$0{,}056$ (A) e $0{,}045$ (B) no Gini, e $0{,}082$ (A) e $0{,}089$ (B) na entropia.

\smallskip
\textbf{(b)} Pela simetria, a impureza ponderada de A é
$\tfrac12I(1/3)+\tfrac12I(2/3)=I(1/3)$, e a de B é
$\tfrac1{12}I(0)+\tfrac{11}{12}I(6/11)=\tfrac{11}{12}I(6/11)$. Como $I(6/11)>0$,
\[
  I\big(\tfrac13\big)<\tfrac{11}{12}\,I\big(\tfrac6{11}\big)
  \iff \frac{I(1/3)}{I(6/11)}<\frac{11}{12}\approx0{,}917 .
\]
As razões são $(4/9)/(60/121)=121/135\approx0{,}896$ para o Gini,
$0{,}918/0{,}994\approx0{,}924$ para a entropia e $(1/3)/(5/11)=11/15\approx0{,}733$ para o erro: só
a entropia passa de $11/12$.

A comparação tem uma leitura simples. A divisão A leva os dois filhos da proporção $1/2$ à
proporção $1/3$; a divisão B torna pura $1/12$ das mudas e deixa as outras numa região quase tão
misturada quanto o nó, porque $6/11\approx0{,}545$ está perto de $1/2$. Com
$I(6/11)\approx I(1/2)$, a pergunta é, em essência, se $I(1/3)$ fica abaixo de $11/12$ do
máximo: $G(1/3)/G(1/2)=8/9\approx0{,}889$, mas $D(1/3)/D(1/2)\approx0{,}918$. A entropia é mais
achatada perto do máximo: trata uma região com um terço de uma espécie como quase tão impura
quanto uma equilibrada, e por isso dá pouco valor à purificação parcial de A e, em comparação,
mais valor à purificação completa de um pedaço pequeno, que é o que B faz.

\smallskip
\textbf{(c)} Com \texttt{min\_samples\_leaf=5}, a divisão B deixa um filho com uma única muda e
não é admitida; entre as duas, só A resta, e as medidas deixam de discordar. A é a divisão
preferível: separa grupos de $6$ mudas com maiorias claras, enquanto B se apoia numa mancha
que aparece numa muda só, purifica uma observação e deixa as outras tão misturadas quanto
antes, o que dificilmente se repetirá em mudas novas.

O exemplo mostra que o ``quase'' é literal: as duas medidas podem, de fato, escolher divisões
diferentes. É uma questão de frequência. Numa enumeração de todos os pares de divisões
possíveis em nós de $10$ a $40$ observações, de todas as composições de classes, as duas medidas
discordam em cerca de $2\%$ dos pares, e a proporção quase não muda quando se exige que todos os
filhos tenham pelo menos $5$ observações. Aqui a restrição desfaz a discordância porque a
divisão B isola uma única muda; em geral, o tamanho mínimo de folha não basta para isso. O que
torna a escolha entre o Gini e a entropia pouco importante na prática é a raridade da
discordância, e não uma restrição como a do slide.
\end{solucao}
